\chapter*{Acrónimos}
\addcontentsline{toc}{chapter}{Acrónimos}

La siguiente lista recoge los acrónimos empleados en la memoria. Cada uno se define también
en su primer uso dentro del cuerpo del texto, con la forma larga seguida de la sigla. En el
resumen y en el \textit{abstract} se han escrito las formas largas, sin siglas.

\begin{description}
    \item[ADR] Registro de decisión de arquitectura (del inglés \textit{architecture decision record}).
    \item[API] Interfaz de programación de aplicaciones (del inglés \textit{application programming interface}).
    \item[CC0] Dedicación al dominio público \textit{Creative Commons Zero}.
    \item[CSRF] Falsificación de petición en sitios cruzados (del inglés \textit{cross-site request forgery}).
    \item[CSS] Hojas de estilo en cascada (del inglés \textit{cascading style sheets}).
    \item[CSV] Valores separados por comas (del inglés \textit{comma-separated values}).
    \item[DLC] Contenido descargable (del inglés \textit{downloadable content}).
    \item[DRF] \textit{Django REST Framework}.
    \item[DSA] Acuerdo de servicios para desarrolladores (del inglés \textit{developer services agreement}).
    \item[DTO] Objeto de transferencia de datos (del inglés \textit{data transfer object}).
    \item[E2E] De extremo a extremo (del inglés \textit{end to end}).
    \item[GSD] \textit{Get Stuff Done}, método de desarrollo dirigido por especificación para agentes.
    \item[HHI] Índice de Herfindahl-Hirschman (del inglés \textit{Herfindahl-Hirschman index}).
    \item[HTTP] Protocolo de transferencia de hipertexto (del inglés \textit{hypertext transfer protocol}).
    \item[IDF] Frecuencia inversa de documento (del inglés \textit{inverse document frequency}).
    \item[IGDB] \textit{Internet Game Database}.
    \item[JSON] Notación de objetos de JavaScript (del inglés \textit{JavaScript object notation}).
    \item[LTS] Soporte a largo plazo (del inglés \textit{long-term support}).
    \item[MAP] Precisión media (del inglés \textit{mean average precision}).
    \item[MMR] Relevancia marginal máxima (del inglés \textit{maximal marginal relevance}).
    \item[nDCG] Ganancia acumulada descontada normalizada (del inglés \textit{normalised discounted cumulative gain}).
    \item[ORM] Mapeo objeto-relacional (del inglés \textit{object-relational mapping}).
    \item[RAWG] Base de datos de videojuegos usada como fuente de contraste (\textit{RAWG Video Games Database}).
    \item[REST] Transferencia de estado representacional (del inglés \textit{representational state transfer}).
    \item[SHA] Algoritmo de resumen seguro (del inglés \textit{secure hash algorithm}).
    \item[SQL] Lenguaje de consulta estructurado (del inglés \textit{structured query language}).
    \item[SSR] Renderizado en el servidor (del inglés \textit{server-side rendering}).
    \item[TFG] Trabajo Fin de Grado.
    \item[URL] Localizador uniforme de recursos (del inglés \textit{uniform resource locator}).
    \item[UUID] Identificador único universal (del inglés \textit{universally unique identifier}).
    \item[WCAG] Pautas de accesibilidad para el contenido web (del inglés \textit{web content accessibility guidelines}).
\end{description}
